\documentclass[a4paper,12pt,oneside]{book}

% ── Encoding e lingua ────────────────────────────────────────────────
\usepackage[utf8]{inputenc}
\usepackage[T1]{fontenc}
\usepackage[italian]{babel}

% ── Geometria della pagina ───────────────────────────────────────────
\usepackage[top=3cm, bottom=3cm, left=3.5cm, right=2.5cm]{geometry}

% ── Font ─────────────────────────────────────────────────────────────
\usepackage{lmodern}

% ── Pacchetti grafici e tabelle ──────────────────────────────────────
\usepackage{graphicx}
\usepackage{float}
\usepackage{booktabs}
\usepackage{longtable}
\usepackage{array}
\usepackage{tabularx}

% ── Codice sorgente ──────────────────────────────────────────────────
\usepackage{listings}
\usepackage{xcolor}

\definecolor{codebg}{RGB}{245,245,245}
\definecolor{codeframe}{RGB}{200,200,200}
\definecolor{codestring}{RGB}{163,21,21}
\definecolor{codecomment}{RGB}{0,128,0}
\definecolor{codekeyword}{RGB}{0,0,255}

\lstset{
  basicstyle=\ttfamily\footnotesize,
  backgroundcolor=\color{codebg},
  frame=single,
  rulecolor=\color{codeframe},
  breaklines=true,
  breakatwhitespace=false,
  tabsize=2,
  showstringspaces=false,
  keywordstyle=\color{codekeyword},
  commentstyle=\color{codecomment}\itshape,
  stringstyle=\color{codestring},
  numbers=none,
  xleftmargin=1em,
  xrightmargin=1em,
  aboveskip=1em,
  belowskip=1em,
}

\lstdefinelanguage{json}{
  morestring=[b]",
  morecomment=[l]{//},
  literate=
    *{:}{{{\color{codekeyword}{:}}}}{1}
    {,}{{{\color{codekeyword}{,}}}}{1}
    {\{}{{{\color{codekeyword}{\{}}}}{1}
    {\}}{{{\color{codekeyword}{\}}}}}{1}
    {[}{{{\color{codekeyword}{[}}}}{1}
    {]}{{{\color{codekeyword}{]}}}}{1},
}

% ── Hyperlink ────────────────────────────────────────────────────────
\usepackage{hyperref}
\hypersetup{
  colorlinks=true,
  linkcolor=black,
  citecolor=black,
  urlcolor=blue,
  pdfauthor={Davide Leone},
  pdftitle={PCB-CTF: Piattaforma Web per la Formazione sulla Sicurezza Hardware},
}

% ── Spaziatura ───────────────────────────────────────────────────────
\usepackage{setspace}
\onehalfspacing

% ── Caption ──────────────────────────────────────────────────────────
\usepackage[font=small,labelfont=bf]{caption}

% ── Placeholder per figure ───────────────────────────────────────────
\newcommand{\figurePlaceholder}[2]{%
  \begin{figure}[H]
    \centering
    \fbox{\parbox{0.8\textwidth}{\centering\vspace{2cm}\textit{Screenshot da inserire}\vspace{2cm}}}
    \caption{#2}
    \label{fig:#1}
  \end{figure}
}

% ══════════════════════════════════════════════════════════════════════
% DOCUMENTO
% ══════════════════════════════════════════════════════════════════════

\begin{document}

% ── Frontespizio ─────────────────────────────────────────────────────
\begin{titlepage}
  \centering
  \vspace*{2cm}

  {\Large\textsc{Università degli Studi di ...}}\\[0.5cm]
  {\large Dipartimento di ...}\\[0.3cm]
  {\large Corso di Laurea in ...}\\[3cm]

  {\LARGE\bfseries PCB-CTF}\\[0.4cm]
  {\Large Piattaforma Web Interattiva per la Formazione\\sulla Sicurezza Hardware in Formato CTF}\\[3cm]

  \begin{minipage}{0.45\textwidth}
    \begin{flushleft}
      \textbf{Relatore}\\
      Prof. ...
    \end{flushleft}
  \end{minipage}
  \hfill
  \begin{minipage}{0.45\textwidth}
    \begin{flushright}
      \textbf{Candidato}\\
      Davide Leone\\
      Matricola: ...
    \end{flushright}
  \end{minipage}

  \vfill
  {\large Anno Accademico 2025/2026}
\end{titlepage}

% ── Indice ───────────────────────────────────────────────────────────
\frontmatter
\tableofcontents

% ── Corpo ────────────────────────────────────────────────────────────
\mainmatter

\input{introduzione}

% Capitolo 1 — Stato dell'arte (da scrivere)
% \input{capitolo-1}

% =======================================================================
% Capitolo 2 — Analisi dei requisiti
% =======================================================================

\chapter{Analisi dei requisiti}
\label{ch:requisiti}

\section{Contesto operativo e attori del sistema}
\label{sec:contesto-attori}

Prima di procedere con l'elicitazione dei requisiti, è necessario definire con chiarezza il contesto d'uso della piattaforma e gli attori coinvolti.

La piattaforma PCB-CTF è concepita per l'uso in ambito accademico e formativo, dove un docente --- o più in generale un autore di contenuti --- predispone un esercizio di analisi hardware che uno o più studenti dovranno risolvere. I due attori principali operano in momenti e modalità distinte:

\begin{itemize}
  \item Lo \textbf{studente} accede al simulatore via browser, visualizza il circuito stampato, utilizza gli strumenti virtuali messi a disposizione e risolve progressivamente gli obiettivi proposti. Il suo percorso è guidato da istruzioni testuali e suggerimenti, e il completamento di ciascuna fase è attestato dalla scoperta di un frammento della flag finale.

  \item L'\textbf{autore} (docente, tutor, o progettista dell'esercizio) accede al pannello di configurazione, dove dispone i componenti sulla fotografia del PCB, definisce i punti di misura, configura il comportamento del terminale e struttura la sequenza didattica in step e obiettivi. Il suo lavoro produce una configurazione JSON che il simulatore consumerà.
\end{itemize}

Non è previsto, nella versione attuale, un terzo attore con ruolo amministrativo: l'autore ha pieno controllo su tutte le funzionalità di configurazione. Questa semplificazione riflette lo scenario d'uso primario, dove il docente opera in autonomia sulla propria istanza della piattaforma.


\section{Requisiti funzionali}
\label{sec:requisiti-funzionali}

I requisiti funzionali sono stati organizzati per area tematica, riflettendo le macro-funzionalità della piattaforma.

\subsection{Simulazione del circuito stampato}
\label{subsec:rf-simulazione-pcb}

\paragraph{RF-01 --- Visualizzazione del PCB.} Il sistema deve consentire la visualizzazione di una fotografia ad alta risoluzione di un circuito stampato, sulla quale vengono sovrapposti elementi interattivi (componenti cliccabili, pin, overlay degli strumenti). L'interfaccia deve adattarsi alle dimensioni della finestra del browser mantenendo le proporzioni degli elementi.

\paragraph{RF-02 --- Identificazione dei componenti.} Lo studente deve poter individuare e selezionare componenti hardware sul PCB attraverso il click su aree predefinite dall'autore. Il sistema deve fornire feedback visivo per distinguere componenti già identificati da quelli ancora da scoprire.

\paragraph{RF-03 --- Strumento di ingrandimento.} Il sistema deve fornire una lente di ingrandimento virtuale che consenta di esaminare i dettagli del PCB con un livello di zoom configurabile. La lente deve poter essere ancorata in una posizione fissa per liberare il cursore.

\subsection{Strumenti di misura}
\label{subsec:rf-strumenti-misura}

\paragraph{RF-04 --- Multimetro digitale.} Il sistema deve simulare un multimetro digitale con due sonde (rossa e nera), modalità tensione e modalità resistenza. Il multimetro deve visualizzare valori realistici con rumore digitale simulato. Le letture devono rispettare le regole fisiche: in modalità tensione con due sonde, il display mostra la differenza di potenziale; in modalità resistenza, la somma delle resistenze.

\paragraph{RF-05 --- Adattatore UART.} Il sistema deve simulare un adattatore USB-to-serial con pin TX, RX e GND trascinabili. Il collegamento deve rispettare la logica del protocollo UART, in particolare la corrispondenza incrociata tra TX e RX. Il sistema deve validare in tempo reale la correttezza dei collegamenti senza rivelare esplicitamente la soluzione.

\paragraph{RF-06 --- Strumento di firmware dump SPI.} Il sistema deve simulare un programmatore SPI con sei sonde corrispondenti ai segnali del bus (VCC, GND, CS, CLK, MOSI, MISO). Ogni sonda deve poter essere collegata solo al pin con il ruolo corrispondente. Al completamento delle connessioni, il sistema deve simulare il trasferimento del firmware con una barra di progresso e offrire il download di un file binario reale.

\paragraph{RF-07 --- Connessione delle sonde.} Tutti gli strumenti con sonde trascinabili devono implementare un meccanismo di aggancio automatico (snap) che faciliti il collegamento ai pin vicini al cursore. Il raggio di aggancio deve essere sufficientemente ampio da rendere l'interazione confortevole ma non così ampio da generare ambiguità quando i pin sono ravvicinati.

\subsection{Terminale simulato}
\label{subsec:rf-terminale}

\paragraph{RF-08 --- Emulazione del terminale.} Il sistema deve fornire un terminale testuale interattivo che simuli un ambiente a riga di comando, senza eseguire effettivamente comandi sul server. Il terminale deve supportare comandi built-in (\texttt{ls}, \texttt{cd}, \texttt{cat}, \texttt{pwd}, \texttt{grep}, \texttt{find}, \texttt{clear}) e comandi personalizzati definiti dall'autore.

\paragraph{RF-09 --- Filesystem simulato.} Il terminale deve operare su un filesystem virtuale configurabile, con supporto per directory, file con contenuto testuale, permessi e proprietari. Il filesystem deve essere navigabile tramite i comandi standard.

\paragraph{RF-10 --- Sequenza di boot.} Il terminale deve poter simulare una sequenza di avvio del dispositivo, con stadi animati che riproducono il flusso di un boot reale (bootloader, kernel, login). Alcuni stadi devono prevedere interazione da parte dell'utente (inserimento credenziali).

\paragraph{RF-11 --- Sistema di flag terminale.} Il terminale deve implementare un meccanismo di scoperta progressiva delle flag: l'esecuzione di determinati comandi sblocca frammenti della flag, e l'interfaccia mostra in tempo reale il progresso raggiunto.

\paragraph{RF-12 --- Tab multipli.} Il terminale deve supportare tab multipli, ciascuno con il proprio contesto di esecuzione (filesystem, percorso corrente, cronologia). Questo consente di simulare scenari realistici in cui lo studente opera simultaneamente sul dispositivo target (via UART) e sulla propria macchina di analisi.

\paragraph{RF-13 --- Vincoli sui comandi.} Il sistema deve consentire all'autore di definire vincoli sull'esecuzione dei comandi: disponibilità limitata a specifiche directory, requisito di permessi, prerequisiti (comandi già eseguiti o flag già scoperte), e restrizione a specifici stadi di boot.

\subsection{Struttura dell'esercizio}
\label{subsec:rf-struttura-esercizio}

\paragraph{RF-14 --- Organizzazione in step e obiettivi.} L'esercizio deve essere strutturato in step sequenziali, ciascuno con uno o più obiettivi. Lo studente affronta gli step in ordine; all'interno di ogni step, gli obiettivi possono essere risolti nella sequenza prevista dall'autore.

\paragraph{RF-15 --- Tipologie di obiettivi.} Il sistema deve supportare almeno cinque tipologie di obiettivi: identificazione di componenti (click), condizioni su pin (valori di misura), connessione UART, interazione con il terminale (scoperta di flag), ed estrazione firmware SPI.

\paragraph{RF-16 --- Flag progressiva.} Il completamento di ogni step e obiettivo deve contribuire alla costruzione progressiva della flag finale dell'esercizio. L'interfaccia deve mostrare in tempo reale la flag parziale, con indicazione chiara dei frammenti ancora da scoprire.

\paragraph{RF-17 --- Dialogo di completamento.} Al termine dell'esercizio, il sistema deve mostrare un dialogo configurabile che può includere: la flag completa con possibilità di copia, un pulsante di redirect verso una piattaforma esterna, e un pulsante per il download di un file.

\subsection{Authoring e configurazione}
\label{subsec:rf-authoring}

\paragraph{RF-18 --- Editor visuale del PCB.} L'autore deve poter posizionare componenti e pin sulla fotografia del PCB tramite interazione diretta (drag-and-drop), con supporto per zoom, rotazione e traslazione della vista.

\paragraph{RF-19 --- Editor del terminale.} Il sistema deve fornire un editor strutturato per la configurazione del terminale: definizione dei comandi (tipo, output, vincoli, effetti collaterali), organizzazione del filesystem ad albero, configurazione della sequenza di boot, e definizione delle parti della flag.

\paragraph{RF-20 --- Anteprima del terminale.} L'autore deve poter testare il terminale configurato direttamente dal pannello di authoring, senza necessità di passare al simulatore.

\paragraph{RF-21 --- Sistema di preset.} Il sistema deve consentire il salvataggio e il caricamento di configurazioni complete (esercizio + terminale) come preset riutilizzabili. Ogni preset deve essere autocontenuto, includendo una copia dell'immagine del PCB.

\paragraph{RF-22 --- Configurazione degli strumenti built-in.} L'autore deve poter configurare i parametri degli strumenti built-in: raggio e livello di zoom della lente, comportamento dell'adattatore UART (persistenza, associazione con il terminale), parametri del firmware dump (sonde, durata, file binario).

\paragraph{RF-23 --- Gruppi di strumenti.} L'autore deve poter definire gruppi di strumenti attivabili contemporaneamente, superando il vincolo predefinito di mutua esclusività.

\subsection{Lancio automatico del terminale}
\label{subsec:rf-auto-lancio}

\paragraph{RF-24 --- Auto-lancio del terminale.} Il terminale deve poter essere attivato automaticamente in risposta a eventi specifici: completamento di una connessione UART, completamento di un dump firmware, o attivazione di un obiettivo di tipo terminale. Il componente terminale da lanciare e il tab iniziale devono essere configurabili.

\subsection{Allegati agli hint}
\label{subsec:rf-hint-files}

\paragraph{RF-25 --- File allegati ai suggerimenti.} L'autore deve poter associare uno o più file scaricabili al suggerimento di un obiettivo. Lo studente, quando richiede il suggerimento, deve poter visualizzare il testo dell'hint e scaricare i file allegati (ad esempio, datasheet, schemi elettrici, immagini di riferimento). I file devono essere caricabili tramite il pannello di authoring e serviti direttamente dal server.


\section{Requisiti non funzionali}
\label{sec:requisiti-non-funzionali}

\subsection{Accessibilità e distribuzione}
\label{subsec:rnf-accessibilita}

\paragraph{RNF-01 --- Zero installazione.} La piattaforma deve funzionare interamente nel browser, senza richiedere l'installazione di software aggiuntivo, plugin o estensioni lato studente. L'unico prerequisito deve essere un browser moderno con supporto JavaScript.

\paragraph{RNF-02 --- Deployment semplificato.} Il sistema deve produrre un artefatto di deployment autocontenuto (standalone), deployabile in ambienti containerizzati (Docker) senza dipendenze runtime esterne. Non deve essere richiesto un database o servizi esterni.

\subsection{Usabilità}
\label{subsec:rnf-usabilita}

\paragraph{RNF-03 --- Guidabilità.} Ogni obiettivo deve essere accompagnato da un'istruzione testuale e da un suggerimento opzionale, in modo che lo studente possa procedere anche senza conoscenze pregresse specifiche sull'hardware in esame.

\paragraph{RNF-04 --- Feedback visivo.} Il sistema deve fornire feedback visivo immediato per tutte le interazioni: evidenziazione dei pin nell'area di snap, animazione dei fili durante il trascinamento, indicazione cromatica dello stato delle connessioni, barra di progresso durante il dump firmware.

\paragraph{RNF-05 --- Responsività.} L'interfaccia deve adattarsi a diverse risoluzioni e dimensioni dello schermo senza compromettere la funzionalità. Il sistema di coordinate normalizzate deve garantire il corretto posizionamento degli elementi interattivi indipendentemente dalla dimensione della finestra.

\subsection{Prestazioni}
\label{subsec:rnf-prestazioni}

\paragraph{RNF-06 --- Esecuzione lato client.} Tutte le interazioni dello studente --- misurazione con il multimetro, connessione delle sonde, esecuzione dei comandi nel terminale --- devono essere elaborate lato client, senza round-trip al server, per garantire latenza minima e responsività immediata.

\paragraph{RNF-07 --- Indipendenza dal server durante la simulazione.} Una volta caricata la configurazione iniziale, il simulatore deve funzionare senza ulteriori chiamate al server. La perdita di connettività non deve interrompere la sessione in corso.

\subsection{Manutenibilità e estensibilità}
\label{subsec:rnf-manutenibilita}

\paragraph{RNF-08 --- Tipizzazione statica.} Il modello dati deve essere integralmente tipizzato con TypeScript, in modo da intercettare errori di struttura in fase di compilazione.

\paragraph{RNF-09 --- Modularità del terminale.} Il sistema terminale deve supportare la definizione di più componenti terminale indipendenti all'interno di un singolo esercizio, per consentire scenari con dispositivi multipli o con ambienti di analisi separati.

\paragraph{RNF-10 --- Configurazione dichiarativa.} Il comportamento del terminale --- comandi, vincoli, output, effetti collaterali --- deve essere definibile in modo dichiarativo tramite configurazione JSON, senza richiedere modifiche al codice sorgente per la creazione di nuovi esercizi.

\subsection{Sicurezza}
\label{subsec:rnf-sicurezza}

\paragraph{RNF-11 --- Nessuna esecuzione server-side di comandi.} Il terminale simulato non deve mai eseguire comandi reali sul sistema operativo del server. Ogni interazione deve essere gestita dal motore di esecuzione lato client, operando esclusivamente sulle strutture dati in memoria.

\paragraph{RNF-12 --- Isolamento delle sessioni.} Lo stato della sessione di uno studente deve essere completamente indipendente e non deve interferire con le sessioni di altri studenti o con la configurazione dell'autore.


\section{Casi d'uso principali}
\label{sec:casi-uso}

Si descrivono di seguito i casi d'uso più rilevanti, organizzati per attore.

\subsection{Casi d'uso dello studente}
\label{subsec:cu-studente}

\paragraph{CU-01 --- Esplorazione del PCB.} Lo studente apre il simulatore, visualizza il circuito stampato e utilizza la lente di ingrandimento per esaminare i dettagli. Identifica i componenti hardware cliccando sulle aree corrispondenti. Al completamento di ogni identificazione, riceve un frammento della flag.

\paragraph{CU-02 --- Misurazione elettrica.} Lo studente seleziona il multimetro dalla sidebar e trascina le sonde sui punti di test del circuito. Il display mostra i valori di tensione o resistenza. L'obiettivo viene completato quando le sonde sono posizionate sui pin che soddisfano le condizioni definite dall'autore.

\paragraph{CU-03 --- Connessione UART.} Lo studente seleziona lo strumento delle sonde UART e collega i pin dell'adattatore ai pin del circuito, rispettando la logica di crossover del protocollo. Al completamento della connessione corretta, il terminale si apre automaticamente simulando una sessione seriale.

\paragraph{CU-04 --- Estrazione firmware.} Lo studente seleziona lo strumento di firmware dump e collega le sonde SPI ai pin corrispondenti sul PCB. Una volta completate le connessioni richieste, avvia il dump e scarica il file firmware al termine del trasferimento simulato.

\paragraph{CU-05 --- Esplorazione del terminale.} Lo studente interagisce con il terminale simulato: esplora il filesystem, esegue comandi, analizza file di configurazione, scopre credenziali e vulnerabilità. Ogni scoperta significativa sblocca un frammento della flag terminale.

\paragraph{CU-06 --- Completamento dell'esercizio.} Lo studente completa tutti gli step dell'esercizio, inserisce la flag finale nel dialogo di validazione, e riceve conferma del completamento. Il dialogo finale può offrire opzioni di redirect o download.

\subsection{Casi d'uso dell'autore}
\label{subsec:cu-autore}

\paragraph{CU-07 --- Configurazione del PCB.} L'autore carica la fotografia del circuito, posiziona i componenti cliccabili e definisce i pin di misura (con valori elettrici), i pin UART (con ruoli) e i pin SPI (con ruoli). Utilizza il canvas interattivo per il posizionamento preciso.

\paragraph{CU-08 --- Definizione dell'esercizio.} L'autore crea step sequenziali, ciascuno con obiettivi tipizzati. Per ogni obiettivo definisce istruzioni, suggerimenti, frammento di flag e condizioni di completamento specifiche per il tipo (click su componente, condizioni su pin, flag terminale, connessione UART, dump firmware).

\paragraph{CU-09 --- Configurazione del terminale.} L'autore crea uno o più componenti terminale, ciascuno con tab, comandi, filesystem, sequenza di boot e flag. Utilizza l'editor strutturato per definire il comportamento di ogni comando e l'anteprima integrata per verificare il risultato.

\paragraph{CU-10 --- Gestione dei preset.} L'autore salva la configurazione corrente come preset, la carica da un preset esistente, o aggiorna un preset precedentemente salvato. Il sistema traccia le modifiche non salvate rispetto al preset attivo.


\section{Vincoli e assunzioni}
\label{sec:vincoli-assunzioni}

La progettazione ha tenuto conto di un insieme di vincoli e assunzioni derivanti dal contesto d'uso:

\paragraph{V-01 --- Singolo utente simultaneo per istanza.} La piattaforma è progettata per un uso single-user per istanza. Non è prevista la gestione di sessioni concorrenti con autenticazione: ogni studente opera sulla propria copia dell'applicazione, o su un'istanza condivisa dove le sessioni sono indipendenti e non persistenti.

\paragraph{V-02 --- Sessioni non persistenti.} Il progresso dello studente nel simulatore non viene salvato tra le sessioni. Ogni caricamento della pagina riparte dallo stato iniziale dell'esercizio. Questa scelta è coerente con l'uso in contesto laboratoriale, dove la sessione ha una durata definita e il completamento avviene tipicamente in un'unica seduta.

\paragraph{V-03 --- Connettività richiesta solo al caricamento.} La connessione di rete è necessaria per il caricamento iniziale dell'applicazione e della configurazione. Una volta caricata, la simulazione opera interamente lato client e non richiede ulteriori comunicazioni con il server.

\paragraph{V-04 --- Browser moderno.} Si assume che gli studenti utilizzino un browser moderno (Chrome, Firefox, Edge, Safari nelle versioni correnti) con supporto completo per ES2017+, CSS Grid, e le API Web standard (\texttt{localStorage}, Fetch, DOM Events).

\paragraph{V-05 --- Conoscenze tecniche dell'autore.} Si assume che l'autore abbia familiarità con i concetti di sicurezza hardware trattati (UART, SPI, analisi firmware) e con la struttura logica di un sistema embedded. Il pannello di configurazione non fornisce spiegazioni didattiche sui protocolli, ma offre un'interfaccia strutturata per la loro configurazione.

\input{capitolo-3}
\input{capitolo-4}
\chapter{Caso di studio}
\label{ch:caso-di-studio}

Per verificare le capacità espressive della piattaforma e illustrarne concretamente l'uso, è stato progettato un esercizio CTF completo basato sull'analisi di un dispositivo reale: il router TP-Link WR841N, un dispositivo consumer ampiamente diffuso e frequentemente utilizzato nella ricerca sulla sicurezza dei dispositivi embedded.

La scelta di questo dispositivo non è casuale. Il WR841N è stato oggetto di numerosi studi e disclosure di vulnerabilità nel corso degli anni, il che lo rende un soggetto ideale per un esercizio didattico: è sufficientemente semplice da essere accessibile a studenti senza esperienza pregressa, ma presenta una superficie di attacco abbastanza ampia da consentire un percorso formativo articolato, che spazia dall'ispezione visiva del circuito stampato fino all'analisi delle vulnerabilità del software.


\section{Progettazione dell'esercizio}
\label{sec:progettazione-esercizio}

\subsection{Struttura dell'esercizio}
\label{subsec:struttura-esercizio}

L'esercizio è articolato in due step formali, corrispondenti alle fasi di analisi hardware, seguiti da una fase di esplorazione terminale che si attiva automaticamente al completamento della connessione UART:

\textbf{Step 1 --- Hardware Analysis.} Lo studente deve identificare visivamente i componenti principali sulla fotografia del PCB del router: il System-on-Chip (SoC) QCA953x e la memoria EPROM. Successivamente, deve collegare le sonde del firmware dumper ai pin della memoria per estrarre il firmware. Questo step introduce le competenze fondamentali di ricognizione fisica e di interazione con i bus di comunicazione.

\textbf{Step 2 --- UART.} Lo studente deve individuare i pin della porta seriale UART sul circuito e collegare l'adattatore USB-to-serial rispettando la logica del protocollo: TX dell'adattatore su RX del dispositivo, RX su TX, GND su GND. Il completamento di questo step --- che attesta la corretta comprensione del protocollo --- innesca automaticamente l'apertura del terminale simulato, grazie alla configurazione \texttt{terminalComponentId} dell'adattatore UART.

\textbf{Fase terminale (auto-lanciata).} Il terminale, avviato automaticamente dopo la connessione UART, riproduce il comportamento del router TP-Link con una sequenza di boot realistica (U-Boot $\rightarrow$ kernel Linux $\rightarrow$ login). Il componente terminale è configurato con un sistema di sei flag progressive che guidano l'analisi:

\begin{enumerate}
    \item \textit{Boot Analysis} --- analisi delle variabili d'ambiente del bootloader U-Boot tramite il comando \texttt{printenv}, alla ricerca di inconsistenze nella configurazione (ad esempio, \texttt{rootfstype=jffs2} nel bootloader mentre il kernel monta un filesystem SquashFS);
    \item \textit{Root Access} --- ottenimento dell'accesso root tramite credenziali di default del dispositivo (\texttt{root} / \texttt{sohoadmin});
    \item \textit{Hash Cracking} --- esame del file \texttt{/etc/shadow} per estrarre gli hash delle password e utilizzo del tab ``Local Machine'' per crackarli con strumenti come \texttt{hashcat};
    \item \textit{Data Leak} --- analisi delle partizioni del dispositivo con il comando \texttt{strings} per estrarre credenziali WiFi e configurazioni in chiaro dalla partizione di configurazione;
    \item \textit{Command Injection} --- analisi del binario del web server (\texttt{/usr/bin/httpd}) per individuare funzioni vulnerabili a injection di comandi;
    \item \textit{Backdoor Discovery} --- identificazione di processi sospetti tramite \texttt{ps} e analisi di binari anomali in \texttt{/usr/bin/}.
\end{enumerate}

Ogni scoperta sblocca un frammento della flag terminale; al completamento di tutti e sei, la flag composta \texttt{flag\{b00t\_r00t\_h4sh\_l34k\_1nj3ct\_sh3ll\}} viene rivelata. Questa architettura --- step formali per l'interazione hardware, fase terminale auto-lanciata per l'analisi software --- sfrutta il meccanismo di auto-lancio descritto nel Capitolo 4 (\S4.2.2) per collegare fluidamente le due fasi dell'esercizio.

\subsection{Configurazione del terminale}
\label{subsec:configurazione-terminale}

Il terminale dell'esercizio è configurato con due tab:

\begin{itemize}
    \item \textbf{UART Console}: simula la connessione seriale al dispositivo. Include una sequenza di boot completa che riproduce fedelmente l'output del bootloader U-Boot del QCA953x (banner, inizializzazione RAM, checksum del kernel) seguito dal boot Linux con i messaggi del kernel. Il filesystem replica la struttura reale del WR841N: directory \texttt{/etc} con \texttt{passwd}, \texttt{shadow}, \texttt{inittab}; \texttt{/proc} con \texttt{cpuinfo}, \texttt{mtd}, \texttt{meminfo}; \texttt{/web/oem/model.conf} con il modello e la versione del firmware.

    \item \textbf{Local Machine}: simula la macchina di analisi dello studente, con strumenti come \texttt{hashcat} e \texttt{john} disponibili per il cracking degli hash estratti dal dispositivo target.
\end{itemize}

\subsection{Configurazione degli strumenti}
\label{subsec:configurazione-strumenti}

L'esercizio utilizza la seguente configurazione degli strumenti:

\begin{itemize}
    \item La \textbf{lente di ingrandimento} è configurata con raggio 90px e zoom 2.4x, parametri calibrati sulla risoluzione della fotografia del PCB utilizzata.
    \item L'\textbf{adattatore UART} è configurato con persistenza post-connessione e con associazione al componente terminale, in modo che il terminale si apra automaticamente dopo il collegamento.
    \item Il \textbf{firmware dumper} utilizza due sonde (GND e VCC) con mapping esatto verso pin specifici, durata del dump di 1 secondo, e un file firmware binario reale (\texttt{fw\_test.bin}) scaricabile dallo studente.
    \item Il \textbf{dialogo di completamento} mostra un messaggio di congratulazioni con la possibilità di copiare la flag.
\end{itemize}


\section{Walkthrough dello studente}
\label{sec:walkthrough-studente}

Si descrive di seguito il percorso tipico di uno studente che affronta l'esercizio, evidenziando le interazioni con la piattaforma in ciascuna fase.

\subsection{Fase 1 --- Ricognizione del PCB}
\label{subsec:fase1-ricognizione}

\figurePlaceholder{5-1}{Schermata educativa del primo step, con la descrizione dell'esercizio e il pulsante per avviare la fase attiva.}

Lo studente accede alla piattaforma e visualizza la schermata educativa del primo step, che introduce il contesto dell'esercizio. Dopo aver letto le istruzioni, avvia lo step e si trova di fronte alla fotografia del PCB del router.

\figurePlaceholder{5-2}{PCB del TP-Link WR841N nel simulatore, con la lente di ingrandimento posizionata sul SoC QCA953x e il componente evidenziato dopo l'identificazione.}

Utilizzando la lente di ingrandimento, esamina il circuito alla ricerca dei componenti descritti nelle istruzioni. Identifica il SoC cliccando sull'area corrispondente --- il sistema evidenzia il componente e rivela il primo frammento della flag. Ripete l'operazione per la memoria EPROM.

Per il terzo obiettivo dello step, seleziona lo strumento firmware dump dalla sidebar. Due sonde appaiono sullo schermo; lo studente le trascina sui pin indicati. Quando entrambe le sonde sono correttamente collegate (il meccanismo di snap facilita il posizionamento), il pulsante ``DUMP FIRMWARE'' si abilita. Avvia il dump, osserva la barra di progresso, e al termine scarica il file binario del firmware. L'ultimo frammento della flag dello step viene rivelato.

\subsection{Fase 2 --- Connessione UART}
\label{subsec:fase2-uart}

\figurePlaceholder{5-3}{Adattatore UART con le sonde collegate ai pin del PCB: la connessione incrociata TX$\rightarrow$RX è evidenziata dai fili colorati.}

Nel secondo step, lo studente deve stabilire la connessione seriale. Seleziona lo strumento UART e trascina le sonde dell'adattatore sui pin del PCB. Se tenta di collegare TX su TX, la connessione viene rifiutata --- un feedback immediato che lo invita a riflettere sulla logica del protocollo. Dopo aver individuato il mapping corretto (TX$\rightarrow$RX, RX$\rightarrow$TX, GND$\rightarrow$GND), la connessione viene stabilita e il terminale si apre automaticamente.

\subsection{Fase 3 --- Esplorazione del terminale}
\label{subsec:fase3-terminale}

\figurePlaceholder{5-4}{Terminale simulato durante la sequenza di boot del WR841N, con i messaggi del bootloader U-Boot e l'inizializzazione del kernel Linux.}

Il terminale mostra la sequenza di boot del dispositivo: il banner U-Boot, l'inizializzazione dell'hardware, il caricamento del kernel. Lo studente osserva il flusso dei messaggi, che replica fedelmente quello di un dispositivo reale.

Arrivato al prompt U-Boot (\texttt{U-Boot> }), esegue il comando \texttt{printenv} per esaminare le variabili d'ambiente del bootloader. Nota che \texttt{rootfstype=jffs2} mentre il kernel successivamente monta un filesystem SquashFS --- un'inconsistenza che suggerisce una modifica non autorizzata. Questo sblocca il primo frammento della flag terminale.

Dopo il boot del kernel, il dispositivo presenta un prompt di login. Lo studente prova le credenziali di default: \texttt{root} come utente e \texttt{sohoadmin} come password. L'accesso viene concesso e il secondo frammento viene sbloccato.

Con l'accesso root, lo studente esamina il file \texttt{/etc/shadow} per estrarre gli hash delle password. Passa al tab ``Local Machine'' e utilizza \texttt{hashcat} per crackare l'hash MD5. Scopre la password in chiaro e sblocca il terzo frammento.

Prosegue con l'analisi delle partizioni: \texttt{strings /dev/mtdblock3} rivela credenziali WiFi e configurazioni in chiaro nella partizione di configurazione --- un classico esempio di data leak in dispositivi embedded. Quarto frammento sbloccato.

L'analisi del web server (\texttt{strings /usr/bin/httpd | grep exec}) rivela la presenza di funzioni vulnerabili a command injection come \texttt{execFormatCmd}. Quinto frammento.

Infine, il comando \texttt{ps} rivela un processo sospetto. L'analisi del binario \texttt{/usr/bin/backdoorTest} con \texttt{strings} e \texttt{file} ne conferma la natura malevola. Sesto e ultimo frammento sbloccato.

La flag terminale è ora completa; il sistema combina i frammenti degli step precedenti con quelli del terminale per costruire la flag finale dell'esercizio. Lo studente la inserisce nel dialogo di validazione e riceve la conferma del completamento.

\figurePlaceholder{5-5}{Terminale con la flag progressiva completata e il dialogo di completamento dell'esercizio con la flag finale.}


\section{Configurazione lato autore}
\label{sec:configurazione-autore}

\subsection{Creazione dell'esercizio}
\label{subsec:creazione-esercizio}

\figurePlaceholder{5-6}{Pannello di authoring con la fotografia del PCB del WR841N e i componenti posizionati tramite drag-and-drop.}

La configurazione dell'esercizio descritto è stata interamente realizzata tramite il pannello di authoring della piattaforma, senza modifiche dirette ai file JSON.

Il processo è partito dal caricamento della fotografia del PCB nella tab ``Immagine'' delle impostazioni. L'autore ha poi posizionato i componenti cliccabili (SoC, EPROM) direttamente sull'immagine tramite drag-and-drop, regolando le dimensioni degli overlay fino a coprire le aree corrette del circuito.

I pin di misura e i pin UART sono stati posizionati in modo analogo, ciascuno con il ruolo e i valori elettrici appropriati. Per il firmware dump, l'autore ha configurato le sonde e caricato il file binario tramite l'apposito pannello.

\subsection{Configurazione del terminale}
\label{subsec:configurazione-terminale-autore}

\figurePlaceholder{5-7}{Editor del terminale con la configurazione dei comandi del tab UART Console e la struttura del filesystem del WR841N.}

La configurazione del terminale ha rappresentato la parte più laboriosa. L'autore ha utilizzato l'editor strutturato per:

\begin{itemize}
    \item definire la struttura del filesystem tab per tab, creando le directory e i file che replicano il sistema del WR841N;
    \item configurare i comandi built-in (\texttt{ls}, \texttt{cd}, \texttt{cat}, \texttt{pwd}, \texttt{grep}, \texttt{find}, \texttt{clear}) con i vincoli di boot stage appropriati --- ad esempio, i comandi del filesystem sono disponibili solo nello stadio \texttt{shell}, non nel prompt U-Boot;
    \item creare i comandi custom (\texttt{printenv}, \texttt{version}, \texttt{strings}, \texttt{ps}, \texttt{hashcat}) con i rispettivi output statici, condizionali o lookup;
    \item definire i side effect di sblocco flag per ciascun comando chiave;
    \item configurare la sequenza di boot con gli stadi U-Boot, kernel, login e shell.
\end{itemize}

L'anteprima integrata ha permesso di testare iterativamente il comportamento del terminale durante la configurazione, verificando che i comandi producano l'output atteso e che le flag si sblocchino correttamente.

\subsection{Salvataggio come preset}
\label{subsec:salvataggio-preset}

Al termine della configurazione, l'esercizio è stato salvato come preset con il nome ``CTF-1'' e la descrizione ``TP-LINK''. Il sistema ha creato automaticamente una copia dell'immagine del PCB dedicata al preset, garantendo che l'esercizio sia riproducibile anche se l'immagine di lavoro viene successivamente modificata.

Il preset così creato può essere caricato in qualsiasi istanza della piattaforma, permettendo al docente di distribuire l'esercizio ai propri studenti con un semplice file JSON, oppure di condividerlo con colleghi che utilizzano la stessa piattaforma.

\chapter{Validazione e discussione}
\label{ch:validazione}

\section{Criteri di valutazione}
\label{sec:criteri-valutazione}

La validazione di una piattaforma didattica presenta sfide metodologiche specifiche: l'efficacia formativa è intrinsecamente difficile da misurare in modo quantitativo, poiché dipende da variabili quali il livello di partenza degli studenti, il contesto di utilizzo e la qualità degli esercizi proposti. In assenza di un campione statisticamente significativo di utenti --- la piattaforma è stata sviluppata nell'ambito di questo lavoro di tesi e non è ancora stata sottoposta a una sperimentazione su larga scala --- la validazione si articola su tre assi qualitativi: la copertura dei concetti didattici, la fedeltà della simulazione rispetto all'esperienza reale, e l'espressività del sistema di authoring.

\subsection{Copertura dei concetti didattici}

Il percorso formativo coperto dalla piattaforma abbraccia le fasi principali dell'analisi hardware di un dispositivo embedded:

\begin{table}[H]
\centering
\caption{Copertura dei concetti didattici nelle fasi dell'analisi hardware}
\begin{tabularx}{\textwidth}{X X X}
\toprule
\textbf{Fase dell'analisi} & \textbf{Strumento/funzionalità} & \textbf{Concetto didattico} \\
\midrule
Ricognizione fisica & Lente di ingrandimento, click su componenti & Identificazione dei chip, lettura delle sigle \\
Misurazioni elettriche & Multimetro (tensione, resistenza) & Individuazione dei livelli logici, test di continuità \\
Identificazione interfacce & Pin UART (TX, RX, GND) & Riconoscimento dei bus di comunicazione seriale \\
Connessione seriale & Adattatore UART con crossover & Protocollo UART, logica TX/RX incrociata \\
Estrazione firmware & Sonde SPI (6 segnali) & Bus SPI, ruoli dei segnali, dump della memoria flash \\
Analisi bootloader & Terminale con boot sequence & U-Boot, variabili d'ambiente, configurazione del boot \\
Accesso al sistema & Login con credenziali & Credenziali di default, hardening insufficiente \\
Analisi del filesystem & Comandi built-in (\texttt{ls}, \texttt{cat}, \texttt{grep}, \texttt{find}) & Struttura di un sistema Linux embedded \\
Cracking di password & Tab ``Local Machine'' con hashcat & Hash MD5, attacchi dizionario, \texttt{/etc/shadow} \\
Analisi delle vulnerabilità & Comandi \texttt{strings}, \texttt{ps} & Command injection, backdoor, data leak \\
\bottomrule
\end{tabularx}
\end{table}

Questa copertura corrisponde al percorso tipico descritto nei testi di riferimento sull'analisi hardware, come ``The Hardware Hacking Handbook'' (Jasper van Woudenberg, Colin O'Flynn) e ``Practical IoT Hacking'' (Fotios Chantzis et al.), coprendo le fasi dalla ricognizione fisica fino all'analisi delle vulnerabilità software.

Un'area non coperta dalla piattaforma nella sua versione attuale è quella degli attacchi side-channel (analisi di consumo energetico, emissioni elettromagnetiche, timing attack) e del protocollo JTAG, che richiederebbero un livello di simulazione significativamente più complesso e che sono tipicamente affrontati in fasi avanzate della formazione.

\subsection{Fedeltà della simulazione}

La simulazione implementata opera necessariamente a un livello di astrazione superiore rispetto all'interazione con hardware reale. Tuttavia, diversi accorgimenti contribuiscono a ridurre il divario percettivo:

\begin{itemize}
    \item Il \textbf{rumore digitale} del multimetro riproduce l'instabilità tipica delle ultime cifre di uno strumento reale, insegnando implicitamente allo studente a interpretare le letture con la corretta incertezza.
    \item La \textbf{validazione basata sui ruoli} delle connessioni UART e SPI obbliga lo studente a ragionare sulla logica dei protocolli, non semplicemente a ``indovinare'' il collegamento corretto.
    \item La \textbf{sequenza di boot} del terminale replica fedelmente l'output di un dispositivo reale, con tempi di animazione calibrati per risultare credibili.
    \item Il \textbf{filesystem} include file e directory coerenti con il sistema operativo del dispositivo simulato, non placeholder generici.
\end{itemize}

Il limite principale riguarda l'assenza di variabilità: in un dispositivo reale, le misurazioni sono soggette a rumore ambientale, i collegamenti possono essere intermittenti, e il comportamento del sistema può variare tra sessioni. La simulazione, per sua natura deterministica, non cattura questa complessità. Questa è una limitazione consapevole, non un difetto: in fase didattica, il determinismo è un vantaggio, poiché consente allo studente di concentrarsi sui concetti senza le frustrazioni legate a problemi di contatto o interferenze.

\subsection{Espressività del sistema di authoring}

Il sistema di authoring è stato progettato per consentire la creazione di esercizi eterogenei senza modifiche al codice sorgente. Per valutarne l'espressività, consideriamo la varietà di scenari configurabili:

\begin{itemize}
    \item \textbf{Esercizi di sola ricognizione}: step con obiettivi di tipo ``component'', dove lo studente identifica i chip utilizzando la lente di ingrandimento. Non richiedono terminale né connessioni.
    \item \textbf{Esercizi di misura elettrica}: step con obiettivi di tipo ``pin'', dove lo studente utilizza il multimetro per individuare livelli logici o verificare continuità.
    \item \textbf{Esercizi di connessione}: step con obiettivi UART o firmware-dump, focalizzati sulla comprensione dei protocolli di comunicazione.
    \item \textbf{Esercizi di analisi terminale}: step con obiettivi terminale complessi, sequenze di boot multi-stadio, flag condizionali.
    \item \textbf{Esercizi combinati}: step che integrano strumenti fisici e terminale, come nell'esercizio del caso di studio.
\end{itemize}

Il sistema di sei tipi di output del terminale (statico, condizionale, template, lookup, dinamico, script) consente di simulare una gamma molto ampia di comportamenti, dalla semplice restituzione di testo fisso fino all'esecuzione di logica arbitraria via funzioni JavaScript.

I vincoli sui comandi (percorso, permessi, prerequisiti, argomenti, stadio di boot) permettono di costruire percorsi didattici guidati dove la disponibilità dei comandi evolve con il progresso dello studente, riducendo il rischio che salti fasi o arrivi a risultati senza aver compreso i passaggi intermedi.


\section{Confronto con soluzioni esistenti}
\label{sec:confronto}

Per contestualizzare il contributo di PCB-CTF, è utile confrontarlo con le piattaforme e gli approcci didattici esistenti nel campo della sicurezza hardware.

\begin{table}[H]
\centering
\caption{Confronto tra PCB-CTF e le soluzioni esistenti}
\begin{tabularx}{\textwidth}{l c c c c}
\toprule
\textbf{Caratteristica} & \textbf{CTFd / PicoCTF} & \textbf{RHme (Riscure)} & \textbf{Remote Labs} & \textbf{PCB-CTF} \\
\midrule
Simulazione hardware & No & No (hardware reale) & Parziale & Sì \\
Interazione con PCB & No & Sì (fisico) & Via webcam & Sì (virtuale) \\
Strumenti simulati & No & No (reali) & Reali via rete & Sì \\
Terminale embedded & No & No & Sì (SSH) & Sì (simulato) \\
Accessibilità browser & Sì & No & Parziale & Sì \\
Costo per studente & Zero & Hardware ($\sim$\,\texteuro\,50+) & Infrastruttura & Zero \\
Scalabilità & Alta & Bassa & Media & Alta \\
Authoring no-code & Parziale & No & No & Sì \\
Determinismo & N/A & No & No & Sì \\
Sicurezza server & N/A & N/A & Rischio & Nessun rischio \\
\bottomrule
\end{tabularx}
\end{table}

\textbf{CTFd e PicoCTF} sono piattaforme CTF general-purpose, eccellenti per sfide software (web, crypto, reverse engineering) ma prive di qualsiasi componente di interazione hardware. Lo studente affronta sfide di analisi firmware come esercizi di reverse engineering puro, senza il contesto dell'hardware che lo ospita.

\textbf{RHme (Riscure Hack me)} è una board hardware dedicata alla formazione sulla sicurezza, che richiede l'acquisto fisico del dispositivo. Offre un'esperienza autentica ma con costi e logistica che ne limitano la scalabilità.

\textbf{I laboratori remoti} (come quelli offerti da alcune università e organizzazioni di formazione) forniscono accesso a hardware reale via rete, ma richiedono infrastruttura dedicata, manutenzione, e presentano limiti di concorrenza e disponibilità.

PCB-CTF si posiziona in una nicchia non coperta dalle soluzioni esistenti: offre un'interazione visiva e strumentale con l'hardware --- assente nelle piattaforme CTF tradizionali --- in un formato completamente virtualizzato che non richiede hardware fisico né infrastruttura server complessa.


\section{Limitazioni}
\label{sec:limitazioni}

Nonostante i risultati raggiunti, la piattaforma presenta limitazioni che è necessario riconoscere.

\textbf{Assenza di esecuzione reale.} Il terminale simula, non esegue. Lo studente non può lanciare comandi arbitrari, ma solo quelli previsti dall'autore. Questo implica che la piattaforma non può replicare l'esperienza esplorativa aperta tipica di un'analisi reale, dove il percorso è spesso non lineare e guidato dall'intuizione. Per mitigare questa limitazione, il sistema di output condizionale e lookup consente all'autore di prevedere un insieme ragionevolmente ampio di input, ma la copertura non potrà mai essere completa.

\textbf{Scenari non coperti.} Come già accennato, i protocolli JTAG, I\textsuperscript{2}C, e CAN non sono attualmente supportati. Gli attacchi side-channel --- analisi di potenza, timing, fault injection --- richiederebbero un paradigma di simulazione fondamentalmente diverso e non rientrano nell'ambito di questo lavoro.

\textbf{Singolo utente.} La piattaforma non supporta sessioni concorrenti con autenticazione, scoring, o classifiche. In un contesto di competizione CTF, queste funzionalità sarebbero essenziali.

\textbf{Persistenza limitata.} Il progresso dello studente non viene salvato tra le sessioni. Questo è coerente con l'uso in laboratorio, ma limita l'adozione in contesti di e-learning dove lo studente potrebbe voler riprendere un esercizio in un momento successivo.

\textbf{Validazione empirica.} La piattaforma non è ancora stata sottoposta a una sperimentazione formale con studenti. La valutazione presentata in questo capitolo è qualitativa e basata sull'analisi delle funzionalità; una validazione quantitativa dell'efficacia didattica richiederebbe uno studio controllato con gruppi di studenti, pre-test e post-test, che costituisce un naturale sviluppo futuro.


\section{Sviluppi futuri}
\label{sec:sviluppi-futuri}

I limiti identificati suggeriscono diverse direzioni di evoluzione:

\textbf{Supporto multi-utente e scoring.} L'aggiunta di un sistema di autenticazione e di una classifica trasformerebbe la piattaforma da strumento didattico individuale a piattaforma CTF competitiva, adatta a eventi e competizioni. Questo richiederebbe l'introduzione di un database per la persistenza degli utenti e dei punteggi, e una logica di timing per la classifica.

\textbf{Integrazione con hardware reale.} Un'evoluzione particolarmente interessante sarebbe la possibilità di alternare fasi simulate e fasi su hardware reale all'interno dello stesso esercizio. Ad esempio, lo studente potrebbe completare la ricognizione e le misurazioni in simulazione, e poi passare a un laboratorio remoto per l'estrazione fisica del firmware. Questo approccio ibrido combinerebbe i vantaggi della simulazione (sicurezza, scalabilità, determinismo) con l'autenticità dell'esperienza su hardware reale.

\textbf{Nuovi protocolli e strumenti.} L'architettura del sistema di pin e strumenti è già predisposta per l'estensione. Protocolli come JTAG e I\textsuperscript{2}C potrebbero essere aggiunti come nuovi tipi di pin e strumenti, riutilizzando il meccanismo di connessione role-based e il sistema di validazione esistente.

\textbf{Terminale con LLM.} Un'evoluzione del terminale potrebbe integrare un modello linguistico per gestire comandi non previsti dall'autore, generando risposte plausibili e coerenti con il contesto simulato. Questo amplierebbe significativamente l'esperienza esplorativa dello studente, pur mantenendo il controllo sull'output dei comandi chiave per la progressione dell'esercizio.

\textbf{Persistenza del progresso e analytics.} Il salvataggio del progresso dello studente, combinato con analytics sull'utilizzo (tempo per obiettivo, comandi tentati, errori più frequenti), fornirebbe dati preziosi sia allo studente --- per riprendere il lavoro --- sia al docente --- per identificare i punti di difficoltà e migliorare gli esercizi.

\chapter*{Conclusioni}
\addcontentsline{toc}{chapter}{Conclusioni}

Il presente lavoro ha affrontato il problema della formazione sulla sicurezza hardware in contesti dove l'accesso a laboratori fisici è limitato, proponendo una piattaforma web che simula l'intero flusso di lavoro dell'analisi di un dispositivo embedded --- dall'ispezione visiva del circuito stampato fino all'estrazione e all'analisi del firmware --- in un ambiente completamente virtualizzato e accessibile da browser.

La piattaforma PCB-CTF è stata progettata attorno a due principi guida: la \textbf{fedeltà dell'interazione}, per garantire che l'esperienza dello studente sia trasferibile al contesto reale, e la \textbf{configurabilità dichiarativa}, per consentire ai docenti di creare esercizi personalizzati senza modificare il codice sorgente.

Sul piano dell'interazione, gli strumenti simulati --- multimetro con rumore digitale, adattatore UART con validazione incrociata, programmatore SPI con sei sonde role-based --- replicano la logica operativa delle controparti reali. Il terminale, pur operando su dati predefiniti, riproduce in modo convincente il comportamento di un sistema embedded, dalla sequenza di boot del bootloader U-Boot fino all'esplorazione del filesystem Linux. Il meccanismo di coordinate normalizzate e snap automatico delle sonde rende l'interazione fluida e intuitiva, senza sacrificare il realismo.

Sul piano dell'authoring, il pannello di configurazione consente di definire esercizi arbitrariamente complessi attraverso un editor visuale. Il sistema di sei tipi di output del terminale, combinato con i vincoli contestuali sui comandi e il meccanismo di flag progressive, offre un livello di espressività sufficiente a replicare scenari di analisi realistici, come dimostrato dal caso di studio basato sul router TP-Link WR841N.

Le scelte architetturali si sono rivelate adeguate alla complessità del sistema. La gestione dello stato tramite quattro store Zustand indipendenti ha garantito una separazione chiara delle responsabilità, senza i costi di complessità di un sistema centralizzato. L'assenza di un database e l'adozione di file JSON per la persistenza hanno mantenuto il progetto leggero e immediatamente deployabile, coerentemente con la natura didattica dell'applicazione.

Il caso di studio ha dimostrato che la piattaforma è in grado di supportare un percorso formativo completo, dalla ricognizione hardware all'analisi delle vulnerabilità software, in un formato guidato e gamificato. Lo studente attraversa le stesse fasi concettuali di un'analisi reale --- identificare, misurare, connettersi, estrarre, esplorare, scoprire --- acquisendo competenze metodologiche che costituiscono una base solida per il successivo approccio all'hardware fisico.

Le limitazioni individuate --- l'assenza di esecuzione reale nel terminale, il mancato supporto di protocolli come JTAG e I\textsuperscript{2}C, la mancanza di multi-utenza e persistenza del progresso --- indicano chiaramente le direzioni di evoluzione futura. In particolare, l'integrazione con laboratori remoti per un approccio ibrido simulazione-hardware, e l'introduzione di un sistema di scoring per contesti competitivi, rappresentano le estensioni più promettenti.

In definitiva, PCB-CTF colma un vuoto nell'offerta formativa sulla sicurezza hardware, offrendo un'alternativa accessibile e scalabile ai laboratori fisici. Non si propone di sostituire l'esperienza su hardware reale, ma di preparare lo studente ad affrontarla con le competenze concettuali e la metodologia già acquisite, riducendo la barriera d'ingresso a un campo che, per la crescente pervasività dei dispositivi IoT, è destinato a diventare sempre più rilevante.


% ── Bibliografia ─────────────────────────────────────────────────────
% \bibliographystyle{plain}
% \bibliography{bibliografia}

% ── Appendici ────────────────────────────────────────────────────────
\appendix
\chapter{Struttura JSON di un esercizio completo}
\label{chap:appendice-a}

Si riporta di seguito lo schema completo della struttura dati \texttt{Exercise}, che rappresenta la configurazione di un esercizio nella piattaforma PCB-CTF. Per ciascun campo viene indicato il tipo, l'obbligatorietà e una descrizione sintetica del ruolo all'interno del sistema.

Lo schema è presentato in formato annotato, dove ogni campo è accompagnato da un commento esplicativo. I valori tra parentesi angolari (\texttt{<...>}) indicano il tipo atteso; i campi contrassegnati con \texttt{?} sono opzionali.

\section{Schema dell'esercizio (\texttt{Exercise})}

\begin{lstlisting}[language=json]
{
  // -- Immagine del PCB --
  "pcbImage": "<string>",
  // Percorso dell'immagine del circuito stampato (relativo a /public).
  // Esempio: "/images/pcb_v2.jpg"

  // -- Flag iniziale --
  "initialFlag": "<string>",
  // Template della flag visualizzato allo studente prima che inizi a
  // scoprire i frammenti. I caratteri '?' vengono progressivamente
  // sostituiti dai flagPart degli obiettivi completati.
  // Esempio: "flag{????????????????}"

  // -- Step dell'esercizio --
  "steps": [
    {
      "id": "<string>",
      // Identificatore univoco dello step.
      // Generato con il pattern: "step-" + timestamp base36 + random.

      "title": "<string>",
      // Titolo dello step, visualizzato nell'intestazione del simulatore.

      "description": "<string>",
      // Testo descrittivo mostrato nella schermata educativa iniziale
      // dello step, prima che lo studente avvii la fase attiva.

      "expectedFlag": "<string>",
      // Flag completa attesa per questo step. Viene confrontata con la
      // concatenazione dei flagPart degli obiettivi per verificare il
      // completamento. Esempio: "flag{CPU-123EPROMDUMP}"

      "availableTools?": ["<Tool>"],
      // Lista degli strumenti disponibili durante lo step.
      // Valori ammessi: "pointer", "magnifier", "multimeter",
      //                 "probes", "firmware-dump".
      // Se omesso, tutti gli strumenti sono disponibili.

      "objectives": [
        {
          "id": "<string>",
          // Identificatore univoco dell'obiettivo.

          "name": "<string>",
          // Nome dell'obiettivo, visualizzato nel pannello laterale
          // e nel tracker di progressione.

          "type?": "<ObjectiveType>",
          // Tipologia dell'obiettivo. Determina il meccanismo di
          // completamento. Valori ammessi:
          //   "component"     -- click su un componente del PCB
          //   "pin"           -- connessione sonde a pin specifici
          //   "uart"          -- connessione UART corretta (crossover)
          //   "terminal"      -- scoperta di flag nel terminale
          //   "firmware-dump" -- dump firmware SPI completato
          // Se omesso, il tipo e' "component" (default implicito).

          "componentId?": "<string>",
          // Per type="component": ID del componente hardware collegato.
          // Il click sull'overlay del componente completa l'obiettivo.

          "instruction": "<string>",
          // Istruzione testuale mostrata allo studente quando
          // l'obiettivo e' attivo.

          "hint": "<string>",
          // Suggerimento opzionale, rivelabile dallo studente.

          "flagPart": "<string>",
          // Frammento di flag sbloccato al completamento dell'obiettivo.
          // La concatenazione ordinata dei flagPart di tutti gli
          // obiettivi di uno step deve corrispondere a expectedFlag.

          "coords": ["<number>", "<number>", "<number>", "<number>"],
          // Coordinate normalizzate [x%, y%, width%, height%] dell'area
          // cliccabile sul PCB. Per obiettivi non spaziali (terminal,
          // uart) si usa [0, 0, 0, 0].

          // -- Campi specifici per type="pin" / type="uart" --

          "pinConditions?": [
            {
              "pinId": "<string>",
              // ID del pin sul PCB che deve essere collegato.

              "terminal": "<string>"
              // Identificatore della sonda che deve collegare il pin.
              // Per UART: "adapter-tx", "adapter-rx", "adapter-gnd"
              // Per firmware-dump: "fw-probe:<probeId>"
              // Per multimetro: "probe1", "probe2"
            }
          ],

          "pinLogic?": "<'AND' | 'OR'>",
          // Operatore logico tra le pinConditions.
          // "AND" (default): tutte le condizioni devono essere
          // soddisfatte.
          // "OR": almeno una condizione e' sufficiente.

          // -- Campi specifici per type="terminal" --

          "bootStageConditions?": [
            {
              "bootStageId": "<string>",
              // ID dello stadio di boot associato (informativo).

              "unlockedFlags": ["<string>"],
              // Lista di ID flag del terminale che devono essere
              // scoperte per completare l'obiettivo. L'obiettivo
              // si completa quando TUTTI i flag elencati sono stati
              // sbloccati tramite i sideEffects dei comandi.

              "hint": "<string>"
              // Suggerimento per lo stadio (informativo).
            }
          ],

          "terminalComponentId?": "<string>",
          // ID del componente terminale da attivare quando questo
          // obiettivo diventa corrente.

          "requiresUart?": "<boolean>",
          // Se true, il terminale non si avvia finche' la connessione
          // UART non e' stata stabilita.

          "terminalPersistent?": "<boolean>",
          // Se true, il terminale non puo' essere chiuso dalla toolbar
          // una volta attivato per questo obiettivo.

          // -- Campi specifici per type="firmware-dump" --

          "customToolId?": "<string>",
          // ID del custom tool di tipo firmware-dump collegato
          // all'obiettivo.

          // -- File allegati all'hint --

          "hintFiles?": ["<string>"]
          // Array di percorsi di file scaricabili associati al
          // suggerimento dell'obiettivo. I file vengono caricati
          // dall'autore tramite POST /api/hints/upload e salvati
          // in /public/uploads/hints/ con prefisso timestamp.
          // Esempio: ["/uploads/hints/hint-1772839791936-datasheet.pdf"]
        }
      ]
    }
  ],

  // -- Componenti hardware cliccabili --
  "components": [
    {
      "id": "<string>",
      "name": "<string>",
      "instruction": "<string>",
      "hint": "<string>",
      "flagPart": "<string>",
      "coords": ["<number>", "<number>", "<number>", "<number>"]
      // Stessa semantica dell'obiettivo: overlay cliccabile sul PCB
      // con coordinate normalizzate in percentuale.
    }
  ],

  // -- Pin di misura --
  "pins": [
    {
      "id": "<string>",
      "valueV": "<number>",     // Tensione in Volt letta dal multimetro
      "valueOhm": "<number>",   // Resistenza in Ohm letta dal multimetro
      "coords": ["<number>", "<number>", "<number>", "<number>"]
    }
  ],

  // -- Pin UART --
  "uartPins": [
    {
      "id": "<string>",
      "role": "<UartRole>",     // "tx" | "rx" | "gnd" | "vcc"
      "label": "<string>",     // Etichetta visiva, es. "TX (3.3V)"
      "coords": ["<number>", "<number>", "<number>", "<number>"]
    }
  ],

  // -- Pin SPI per firmware dump (opzionale) --
  "firmwareDumpPins?": [
    {
      "id": "<string>",
      "role": "<SpiRole>",     // "vcc" | "gnd" | "cs" | "clk" | "mosi" | "miso"
      "label": "<string>",
      "coords": ["<number>", "<number>", "<number>", "<number>"]
    }
  ],


  // -- Percorso firmware caricato (opzionale) --
  "firmwarePath?": "<string>",
  // Percorso del file firmware sul server, usato dal firmware dumper
  // built-in. Esempio: "/uploads/firmware-1772839791936-fw_test.bin"

  // -- Gruppi di strumenti (opzionale) --
  "toolGroups?": [
    {
      "id": "<string>",
      "name": "<string>",
      "toolIds": ["<string>"]
      // ID degli strumenti che possono essere attivi contemporaneamente.
      // Gli strumenti nello stesso gruppo coesistono; strumenti non in
      // alcun gruppo sono mutuamente esclusivi.
    }
  ],

  // -- Configurazione strumenti built-in (opzionale) --
  "toolConfig?": {

    "magnifier?": {
      "defaultRadius": "<number>",      // Raggio in pixel della lente
      "defaultZoomLevel": "<number>"    // Fattore di zoom (es. 2.4)
    },

    "uartConnector?": {
      "persistAfterConnection": "<boolean>",
      // Se true, il connettore resta visibile dopo la connessione.

      "visibleInSteps": ["<string>"],
      // ID degli step in cui il connettore resta visibile.

      "terminalComponentId?": "<string>",
      // Terminale da avviare automaticamente dopo connessione UART.

      "terminalDefaultTab?": "<string>"
      // Tab del terminale da attivare (es. "uart").
    },

    "terminal?": {
      "requiresUart": "<boolean>",
      "persistent": "<boolean>",
      "bootStageConditions": [
        {
          "bootStageId": "<string>",
          "unlockedFlags": ["<string>"],
          "hint": "<string>"
        }
      ]
    },

    "firmwareDump?": {
      "probes": [
        {
          "id": "<string>",
          "role": "<SpiRole>",      // "vcc" | "gnd" | "cs" | "clk" | "mosi" | "miso"
          "label": "<string>",
          "color": "<string>"       // Colore hex del filo nel simulatore
        }
      ],
      "requiredConnections?": [
        { "probeId": "<string>", "pinId": "<string>" }
      ],
      // Mapping esplicito sonda->pin. Deprecato: la validazione
      // avviene ora per ruolo (probe.role == pin.role).

      "filePath": "<string>",          // Percorso del file firmware
      "fileName": "<string>",          // Nome file mostrato nel download
      "dumpDurationSec": "<number>",   // Durata progress bar in secondi

      "terminalComponentId?": "<string>",
      // Terminale da avviare dopo il completamento del dump.

      "terminalDefaultTab?": "<string>"
    }
  },

  // -- Dialogo di completamento (opzionale) --
  "completionDialog?": {
    "title": "<string>",
    "description": "<string>",

    "redirectUrl": "<string>",          // URL per pulsante di redirect
    "redirectLabel": "<string>",        // Etichetta del pulsante

    "downloadFilePath": "<string>",     // Percorso file scaricabile
    "downloadLabel": "<string>",        // Etichetta pulsante download
    "downloadFileName": "<string>",     // Nome file per l'utente

    "showCopyFlag": "<boolean>"         // Mostra pulsante copia flag
  }
}
\end{lstlisting}


\section{Tipi enumerati}

Per completezza, si riportano i tipi enumerati utilizzati nello schema.

\subsection{ObjectiveType}

\begin{table}[H]
\centering
\begin{tabularx}{\textwidth}{lX}
\toprule
\textbf{Valore} & \textbf{Descrizione} \\
\midrule
\texttt{component} & Completamento tramite click su un componente del PCB \\
\texttt{pin} & Completamento tramite connessione di sonde a pin specifici \\
\texttt{uart} & Completamento tramite connessione UART corretta (crossover TX/RX) \\
\texttt{terminal} & Completamento tramite scoperta di flag nel terminale simulato \\
\texttt{firmware-dump} & Completamento tramite dump firmware SPI \\
\bottomrule
\end{tabularx}
\caption{Valori ammessi per \texttt{ObjectiveType}.}
\label{tab:objective-type}
\end{table}

\subsection{Tool}

\begin{table}[H]
\centering
\begin{tabularx}{\textwidth}{lX}
\toprule
\textbf{Valore} & \textbf{Descrizione} \\
\midrule
\texttt{pointer} & Cursore standard per interazione con il PCB \\
\texttt{magnifier} & Lente di ingrandimento per ispezione visiva \\
\texttt{multimeter} & Multimetro digitale (tensione e resistenza) \\
\texttt{probes} & Adattatore UART (TX, RX, GND) \\
\texttt{firmware-dump} & Programmatore SPI per estrazione firmware \\
\bottomrule
\end{tabularx}
\caption{Valori ammessi per \texttt{Tool}.}
\label{tab:tool}
\end{table}

\subsection{UartRole}

\begin{table}[H]
\centering
\begin{tabular}{llll}
\toprule
\textbf{Valore} & \textbf{Tensione tipica} & \textbf{Resistenza tipica} \\
\midrule
\texttt{tx} & 3.3 V & 50 $\Omega$ \\
\texttt{rx} & 3.3 V & 10 k$\Omega$ \\
\texttt{gnd} & 0 V & 0 $\Omega$ \\
\texttt{vcc} & 5.0 V & 0 $\Omega$ \\
\bottomrule
\end{tabular}
\caption{Valori ammessi per \texttt{UartRole} con tensione e resistenza tipiche.}
\label{tab:uart-role}
\end{table}

\subsection{SpiRole}

\begin{table}[H]
\centering
\begin{tabular}{llll}
\toprule
\textbf{Valore} & \textbf{Descrizione} & \textbf{Colore nel simulatore} \\
\midrule
\texttt{vcc} & Alimentazione & \texttt{\#EF4444} (rosso) \\
\texttt{gnd} & Massa & \texttt{\#6B7280} (grigio) \\
\texttt{cs} & Chip Select & \texttt{\#A855F7} (viola) \\
\texttt{clk} & Clock & \texttt{\#3B82F6} (blu) \\
\texttt{mosi} & Master Out, Slave In & \texttt{\#F59E0B} (ambra) \\
\texttt{miso} & Master In, Slave Out & \texttt{\#10B981} (verde) \\
\bottomrule
\end{tabular}
\caption{Valori ammessi per \texttt{SpiRole} con colori associati nel simulatore.}
\label{tab:spi-role}
\end{table}


\section{Sistema di coordinate}

Tutte le coordinate nel JSON dell'esercizio utilizzano il \textbf{sistema normalizzato a percentuale} descritto nel Capitolo 3 (\S 3.3.3). Il formato è un array di quattro numeri:

\begin{lstlisting}
[x%, y%, width%, height%]
\end{lstlisting}

\noindent dove \texttt{x\%} e \texttt{y\%} rappresentano la posizione dell'angolo superiore sinistro dell'area rispetto alle dimensioni del contenitore, e \texttt{width\%} e \texttt{height\%} ne definiscono le dimensioni. La conversione in pixel avviene a runtime nel componente \texttt{PCBViewer} in base alle dimensioni effettive del contenitore.

Per obiettivi non spaziali (tipo \texttt{terminal}, \texttt{uart}), le coordinate sono convenzionalmente impostate a \texttt{[0, 0, 0, 0]}.


\section{Esempio completo}

Si riporta un estratto semplificato della configurazione dell'esercizio del caso di studio (Capitolo 5), con due step: identificazione dei componenti e connessione UART.

\begin{lstlisting}[language=json]
{
  "pcbImage": "/images/preset-ctf-1.png",
  "initialFlag": "flag{????????????????}",
  "steps": [
    {
      "id": "step-1",
      "title": "Hardware Analysis",
      "description": "Identificare i componenti principali sulla board...",
      "expectedFlag": "flag{CPU-123EPROMDUMP}",
      "availableTools": ["pointer", "magnifier", "multimeter", "firmware-dump"],
      "objectives": [
        {
          "id": "obj-cpu",
          "name": "CPU",
          "type": "component",
          "componentId": "comp-cpu",
          "instruction": "Identifica il System-on-Chip QCA953x...",
          "hint": "Ha una forma quadrata, e' il chip piu' grande.",
          "flagPart": "CPU-123",
          "coords": [32.07, 49.91, 10.98, 16.11]
        },
        {
          "id": "obj-eprom",
          "name": "EPROM",
          "type": "component",
          "componentId": "comp-eprom",
          "instruction": "Individua la memoria EPROM...",
          "hint": "Ha una forma rettangolare.",
          "flagPart": "EPROM",
          "coords": [60.43, 51.07, 11.11, 28.24]
        },
        {
          "id": "obj-dump",
          "name": "Firmware Dump",
          "type": "pin",
          "instruction": "Collega le sonde del firmware dumper...",
          "hint": "GND e VCC devono corrispondere ai pin indicati.",
          "flagPart": "DUMP",
          "coords": [0, 0, 0, 0],
          "pinConditions": [
            { "pinId": "pin-m-2", "terminal": "fw-probe:fw-probe-vcc" },
            { "pinId": "pin-m-1", "terminal": "fw-probe:fw-probe-gnd" }
          ],
          "pinLogic": "AND"
        }
      ]
    },
    {
      "id": "step-2",
      "title": "UART Connection",
      "description": "Collegare l'adattatore UART ai pin corretti...",
      "expectedFlag": "flag{PIN-M-1,_PIN-M-2,_PIN-M-3}",
      "availableTools": ["magnifier", "multimeter", "probes"],
      "objectives": [
        {
          "id": "obj-uart",
          "name": "Connessione UART",
          "type": "pin",
          "instruction": "Collega TX->RX, RX->TX, GND->GND...",
          "hint": "Ricorda il crossover.",
          "flagPart": "PIN-M-1,_PIN-M-2,_PIN-M-3",
          "coords": [0, 0, 0, 0],
          "pinConditions": [
            { "pinId": "pin-m-1", "terminal": "adapter-rx" },
            { "pinId": "pin-m-2", "terminal": "adapter-tx" },
            { "pinId": "pin-m-3", "terminal": "adapter-gnd" }
          ],
          "pinLogic": "AND"
        }
      ]
    }
  ],
  "components": [
    {
      "id": "comp-cpu",
      "name": "CPU",
      "instruction": "",
      "hint": "",
      "flagPart": "CPU",
      "coords": [32.07, 49.91, 10.98, 16.11]
    },
    {
      "id": "comp-eprom",
      "name": "EPROM",
      "instruction": "",
      "hint": "",
      "flagPart": "EPROM",
      "coords": [60.43, 51.07, 11.11, 28.24]
    }
  ],
  "pins": [
    { "id": "pin-m-1", "valueV": 3.3, "valueOhm": 89,  "coords": [3.06, 42.93, 2, 2] },
    { "id": "pin-m-2", "valueV": 0,   "valueOhm": 0,   "coords": [2.94, 46.25, 2, 2] },
    { "id": "pin-m-3", "valueV": 0,   "valueOhm": 0,   "coords": [2.94, 50.07, 2, 2] }
  ],
  "uartPins": [],
  "toolConfig": {
    "magnifier": { "defaultRadius": 90, "defaultZoomLevel": 2.4 },
    "uartConnector": {
      "persistAfterConnection": true,
      "visibleInSteps": ["step-2"],
      "terminalComponentId": "tc-terminal-1"
    },
    "firmwareDump": {
      "probes": [
        { "id": "fw-probe-gnd", "role": "gnd", "label": "GND", "color": "#6B7280" },
        { "id": "fw-probe-vcc", "role": "vcc", "label": "VCC", "color": "#EF4444" }
      ],
      "filePath": "/uploads/firmware-fw_test.bin",
      "fileName": "fw_test.bin",
      "dumpDurationSec": 1,
      "terminalComponentId": "tc-terminal-1"
    }
  },
  "completionDialog": {
    "title": "Exercise Completed!",
    "description": "Congratulations, you have successfully completed all objectives.",
    "redirectUrl": "",
    "redirectLabel": "Next Exercise",
    "downloadFilePath": "",
    "downloadLabel": "Download File",
    "downloadFileName": "",
    "showCopyFlag": true
  }
}
\end{lstlisting}

\chapter{Schema della configurazione terminale}
\label{chap:appendice-b}

Si riporta lo schema completo della struttura dati \texttt{TerminalConfig}, che definisce il comportamento del terminale simulato nella piattaforma PCB-CTF. Lo schema governa l'intero sistema: tab, comandi, filesystem, sequenza di boot, flag e vincoli di esecuzione.

Nel formato di persistenza, la configurazione terminale è incapsulata in un array di \textbf{componenti terminale}, ciascuno con un proprio identificatore, nome e configurazione indipendente:

\begin{lstlisting}[language=json]
[
  {
    "id": "<string>",           // Identificatore univoco del componente
    "name": "<string>",         // Nome descrittivo (es. "UART Console")
    "config": { /* TerminalConfig */ }
  }
]
\end{lstlisting}

Questa struttura modulare consente di definire più terminali indipendenti nello stesso esercizio (ad esempio, una console UART e una macchina locale), ciascuno con il proprio filesystem, comandi e sequenza di boot.

%---

\section{Struttura principale (\texttt{TerminalConfig})}

\begin{lstlisting}[language=json]
{
  "metadata": {
    "version": "<string>",          // Versione dello schema (es. "1.0.0")
    "name": "<string>",             // Nome della configurazione
    "description": "<string>",      // Descrizione testuale
    "author?": "<string>"           // Autore della configurazione
  },

  "tabs": [ /* TabConfig[] -- vedi B.2 */ ],

  "flags?": { /* FlagSystem -- vedi B.5 */ },

  "globalCommands?": {
    // Comandi disponibili in tutti i tab. Stessa struttura di
    // CommandDefinition (B.3). Utile per comandi trasversali
    // come "help" o "exit".
    "<commandName>": { /* CommandDefinition */ }
  },

  "customFunctions?": {
    // Funzioni JavaScript custom referenziabili dai generatori
    // di output dinamico e dagli effetti collaterali.
    "<functionName>": "<string>"    // Codice della funzione
  }
}
\end{lstlisting}


\section{Configurazione del tab (\texttt{TabConfig})}

Ogni tab rappresenta un ambiente terminale indipendente, con il proprio filesystem, set di comandi, variabili d'ambiente e sequenza di boot opzionale.

\begin{lstlisting}[language=json]
{
  "id": "<string>",                     // Identificatore univoco del tab
  "name": "<string>",                   // Nome visualizzato (es. "UART Console")
  "initialPath?": "<string>",           // Directory iniziale (default: "/")

  "filesystem": { /* FilesystemStructure -- vedi B.4 */ },

  "commands": {
    // Dizionario dei comandi disponibili nel tab.
    // La chiave e' il nome del comando; il valore e' un CommandDefinition.
    "<commandName>": { /* CommandDefinition -- vedi B.3 */ }
  },

  "bootSequence?": {
    // Sequenza di boot animata, riprodotta all'apertura del terminale.
    "initialStage": "<string>",         // ID del primo stadio
    "stages": [
      {
        "id": "<string>",              // Identificatore dello stadio
        "name": "<string>",            // Nome descrittivo
        "lines": ["<string>"],         // Righe di testo animate
        "duration?": "<number>",       // Durata in millisecondi
        "nextStage?": "<string>",      // Stadio successivo (auto-transizione)
        "prompt?": "<string>",         // Prompt dello stadio (se interattivo)
        "autoProgressTimeout?": "<number>"
        // Timeout in millisecondi per l'auto-progressione.
        // Se configurato, lo stadio avanza automaticamente allo
        // stadio successivo dopo il periodo indicato di inattività.
        // L'input dell'utente annulla il timeout in corso.
        // Utile per replicare il countdown "Hit any key to stop
        // autoboot" dei bootloader.
      }
    ]
  },

  "environment?": {
    // Variabili d'ambiente del tab. Accessibili dai comandi template.
    "USER": "<string>",
    "HOME": "<string>",
    "PATH": "<string>",
    "SHELL": "<string>"
    // ... qualsiasi coppia chiave-valore
  },

  "defaultConstraints?": { /* CommandConstraints -- vedi B.3.1 */ }
  // Vincoli applicati a tutti i comandi del tab.
  // I vincoli specificati a livello di singolo comando hanno
  // precedenza su quelli di default.
}
\end{lstlisting}

\subsection{Stadi speciali della sequenza di boot}

Il terminale riconosce tre identificatori di stadio con comportamento hardcoded:

\begin{table}[H]
\centering
\begin{tabularx}{\textwidth}{lX}
\toprule
\textbf{ID stadio} & \textbf{Comportamento} \\
\midrule
\texttt{uboot\_wait} & Mostra un messaggio ``Hit any key to stop autoboot'' con countdown. Se lo studente preme un tasto, transita allo stadio successivo (tipicamente \texttt{uboot\_shell}); altrimenti, il boot prosegue automaticamente. \\
\texttt{login} & Mostra un prompt di login. L'input dello studente viene trattato come nome utente e inoltrato al sistema di comandi per la validazione. \\
\texttt{password} & Mostra un prompt di password (input mascherato). L'input viene inoltrato al sistema di comandi; il comando corrispondente alla password corretta deve produrre la transizione allo stadio \texttt{shell}. \\
\bottomrule
\end{tabularx}
\end{table}


\section{Definizione di un comando (\texttt{CommandDefinition})}

\begin{lstlisting}[language=json]
{
  "name?": "<string>",
  // Nome del comando. Opzionale: se omesso, viene derivato dalla
  // chiave nel dizionario commands.

  "aliases?": ["<string>"],
  // Nomi alternativi che invocano lo stesso comando.
  // Es. ["help"] come alias di "?".

  "description": "<string>",
  // Descrizione del comando (usata internamente e nell'eventuale help).

  "usage?": "<string>",
  // Stringa di utilizzo (es. "cat <file>").

  "handler": "<'builtin' | 'custom' | 'script' | 'dynamic'>",
  // Tipo di handler che gestisce l'esecuzione:
  //   "builtin"  -- implementazione nativa in JavaScript (ls, cd, cat, etc.)
  //   "custom"   -- output definito nella configurazione
  //   "script"   -- esecuzione di una funzione JavaScript custom
  //   "dynamic"  -- generazione dinamica tramite funzione registrata

  "builtinType?": "<string>",
  // Per handler="builtin": nome dell'implementazione nativa.
  // Valori supportati: "ls", "cd", "cat", "pwd", "grep", "find", "clear"

  "constraints?": { /* CommandConstraints -- vedi B.3.1 */ },

  "output?": { /* CommandOutput -- vedi B.3.2 */ },

  "sideEffects?": {
    // Effetti collaterali eseguiti dopo la produzione dell'output.

    "unlockFlags?": [
      // Flag da sbloccare. Puo' essere:
      //   - una stringa (ID del flag, sbloccato incondizionatamente)
      //   - un oggetto FlagUnlock con condizione

      "<string>",
      // oppure:
      {
        "id": "<string>",              // ID del flag da sbloccare
        "condition?": { /* ConditionCheck -- vedi B.3.3 */ }
        // Se presente, il flag viene sbloccato solo se la condizione
        // e' soddisfatta (es. argomento specifico passato al comando).
      }
    ],

    "setState?": {
      // Modifica variabili di stato del terminale.
      // Uso tipico: transizione dello stadio di boot.
      "bootStage": "<string>",
      // ... qualsiasi coppia chiave-valore
    },

    "executeCommand?": "<string>",
    // Comando da eseguire automaticamente dopo quello corrente.

    "customEffect?": "<string>"
    // Nome di una funzione custom da invocare.
  },

  "help?": "<string>"
  // Testo di aiuto esteso.
}
\end{lstlisting}

\subsection{Vincoli del comando (\texttt{CommandConstraints})}

I vincoli determinano le condizioni sotto le quali un comando può essere eseguito. Se un vincolo non è soddisfatto, il comando produce un messaggio di errore configurabile.

\begin{lstlisting}[language=json]
{
  "path?": {
    // Vincolo sul percorso corrente del filesystem.
    "type": "<'exact' | 'contains' | 'startsWith' | 'regex'>",
    "value": "<string | string[]>",     // Valore o lista di valori ammessi
    "errorMessage?": "<string>"
  },

  "permissions?": {
    // Vincolo sui permessi dell'utente.
    "requireRoot?": "<boolean>",        // Richiede utente root
    "requireUser?": "<string>",         // Richiede un utente specifico
    "errorMessage?": "<string>"
  },

  "prerequisites?": {
    // Vincolo su prerequisiti (comandi eseguiti, file esistenti, flag).
    "commands?": ["<string>"],          // Comandi che devono essere stati eseguiti
    "files?": ["<string>"],             // File che devono esistere
    "flags?": ["<string>"],             // Flag che devono essere state scoperte
    "condition?": "<'all' | 'any'>",    // Logica di combinazione (default: "all")
    "errorMessage?": "<string>"
  },

  "arguments?": {
    // Vincolo sugli argomenti passati al comando.
    "min?": "<number>",                 // Numero minimo di argomenti
    "max?": "<number>",                 // Numero massimo di argomenti
    "required?": [
      {
        "index": "<number>",            // Posizione dell'argomento (0-based)
        "pattern?": "<string>",         // Pattern regex che deve matchare
        "errorMessage?": "<string>"
      }
    ],
    "errorMessage?": "<string>"
  },

  "state?": {
    // Vincolo sullo stato del terminale.
    "bootStage?": "<string | string[]>",
    // Stadio (o stadi) di boot in cui il comando e' disponibile.
    // Uso tipico: limitare comandi U-Boot allo stadio "uboot_shell",
    // comandi filesystem allo stadio "shell".

    "variables?": {
      // Variabili di stato che devono avere valori specifici.
      "<key>": "<any>"
    },
    "errorMessage?": "<string>"
  }
}
\end{lstlisting}

\subsection{Tipi di output (\texttt{CommandOutput})}

Il sistema supporta sei tipi di output, ciascuno con semantica e struttura propria.

\subsubsection{Output statico}

\begin{lstlisting}[language=json]
{
  "type": "static",
  "lines": ["<string>"]
  // Array di righe restituite verbatim.
}
\end{lstlisting}

\subsubsection{Output condizionale}

\begin{lstlisting}[language=json]
{
  "type": "conditional",
  "conditions": [
    {
      "if": { /* ConditionCheck */ },
      "then": { /* CommandOutput */ },
      "else?": { /* CommandOutput */ }
    }
  ],
  "default?": { /* CommandOutput */ }
  // Valutazione sequenziale: la prima condizione soddisfatta
  // determina l'output. Se nessuna e' vera, si usa default.
}
\end{lstlisting}

\subsubsection{Output template}

\begin{lstlisting}[language=json]
{
  "type": "template",
  "template": "<string>",
  // Stringa con placeholder {{variabile}}.

  "variables": {
    "<name>": {
      "type": "<'static' | 'computed' | 'random' | 'state'>",
      "value?": "<any>",               // Per type="static"
      "min?": "<number>",              // Per type="random": limite inferiore
      "max?": "<number>",              // Per type="random": limite superiore
      "options?": ["<string>"],        // Per type="random": scelta tra opzioni
      "compute?": "<string>"           // Per type="computed": nome funzione
    }
  }
}
\end{lstlisting}

\subsubsection{Output lookup}

\begin{lstlisting}[language=json]
{
  "type": "lookup",
  "argIndex": "<number>",
  // Indice dell'argomento da confrontare (0-based).

  "matchType": "<'contains' | 'equals' | 'regex'>",
  // Strategia di matching.

  "table": {
    // Mappa valore -> output. La chiave e' il valore da matchare,
    // il valore e' l'array di righe da restituire.
    "<matchValue>": ["<string>"]
  },

  "default?": { /* CommandOutput */ }
  // Output di fallback se nessun match.
}
\end{lstlisting}

\subsubsection{Output dinamico}

\begin{lstlisting}[language=json]
{
  "type": "dynamic",
  "generator": "<string>",     // Nome della funzione generatore
  "args?": ["<any>"]           // Argomenti opzionali
}
\end{lstlisting}

\subsubsection{Output script}

\begin{lstlisting}[language=json]
{
  "type": "script",
  "script": "<string>"         // Nome funzione o script inline
}
\end{lstlisting}

\subsection{Controllo delle condizioni (\texttt{ConditionCheck})}

Le condizioni sono utilizzate sia negli output condizionali sia nello sblocco condizionale delle flag.

\begin{lstlisting}[language=json]
{
  "type": "<'argument' | 'path' | 'file' | 'flag' | 'state' | 'custom'>",

  // Per type="argument":
  "index?": "<number>",            // Indice dell'argomento (0-based)
  "contains?": "<string>",         // L'argomento contiene la stringa
  "equals?": "<string>",           // L'argomento e' esattamente uguale
  "regex?": "<string>",            // L'argomento matcha il pattern regex

  // Per type="file":
  "exists?": "<boolean>",          // Il file esiste/non esiste

  // Per type="state":
  "variable?": "<string>",         // Nome della variabile di stato
  "value?": "<any>",               // Valore atteso

  // Per type="custom":
  "check?": "<string>"             // Nome della funzione di verifica
}
\end{lstlisting}


\section{Struttura del filesystem (\texttt{FilesystemStructure})}

Il filesystem del terminale supporta due formati di definizione: un formato ad \textbf{albero annidato} (più leggibile per l'autore) e un formato \textbf{piatto} (utilizzato internamente dal sistema). Il \texttt{TerminalConfigLoader} converte automaticamente il formato annidato in formato piatto al caricamento.

\begin{lstlisting}[language=json]
{
  "tree?": {
    // Formato annidato: le chiavi sono nomi di file/directory.
    // - Se il valore e' una stringa -> file con quel contenuto
    // - Se il valore e' un oggetto -> directory (ricorsiva)
    // - Se il valore e' un oggetto vuoto {} -> directory vuota

    "etc": {                                    // directory
      "passwd": "root:x:0:0:root:/root:/bin/sh",   // file
      "rc.d": {                                 // sotto-directory
        "rcS": "#!/bin/sh\nmount -a\n..."       // file
      }
    },
    "tmp": {}                                   // directory vuota
  },

  "directories": {
    // Formato piatto: percorso -> lista di entry contenute.
    // Popolato automaticamente dal loader a partire da tree.
    "/": ["etc", "tmp"],
    "/etc": ["passwd", "rc.d"],
    "/etc/rc.d": ["rcS"]
  },

  "files": {
    // Formato piatto: percorso completo -> contenuto o metadati.
    // Puo' essere una stringa (shorthand) o un oggetto FileNode.

    "/etc/passwd": "root:x:0:0:root:/root:/bin/sh",
    // oppure:
    "/etc/passwd": {
      "name": "passwd",
      "type": "file",
      "content": "root:x:0:0:root:/root:/bin/sh",
      "permissions?": "-rw-r--r--",
      "owner?": "root",
      "size?": 37,
      "modified?": "Jun 16 2015",
      "metadata?": {}
    }
  },

  "fileDefaults?": {
    // Valori di default per permessi e proprietario, applicati
    // a tutti i file che non specificano valori propri.
    "permissions?": "<string>",     // Es. "-rw-r--r--"
    "owner?": "<string>"            // Es. "root"
  }
}
\end{lstlisting}

\subsection{Processo di normalizzazione}

Il \texttt{TerminalConfigLoader} esegue la normalizzazione del filesystem in quattro fasi:

\begin{enumerate}
    \item \textbf{Flatten} --- converte il \texttt{tree} annidato in strutture piatte \texttt{directories} e \texttt{files};
    \item \textbf{Normalize} --- assicura che tutte le directory intermedie siano presenti e che la root \texttt{/} esista;
    \item \textbf{Compute} --- calcola le dimensioni dei file e verifica la coerenza dei percorsi;
    \item \textbf{Build} --- genera la struttura finale pronta per l'uso da parte dei comandi built-in.
\end{enumerate}


\section{Sistema di flag (\texttt{FlagSystem})}

\begin{lstlisting}[language=json]
{
  "parts": [
    {
      "id": "<string>",
      // Identificatore univoco del frammento di flag.
      // Referenziato nei sideEffects.unlockFlags dei comandi.

      "part": "<string>",
      // Testo del frammento. Contribuisce alla costruzione
      // progressiva della flag visibile allo studente.

      "description": "<string>",
      // Descrizione del frammento (visibile nel pannello di
      // progressione).

      "hint": "<string>"
      // Suggerimento per lo studente.
    }
  ],

  "completeFlag": "<string>",
  // Flag finale completa, risultante dalla concatenazione di
  // tutti i frammenti. Esempio: "flag{b00t_r00t_h4sh_l34k_1nj3ct_sh3ll}"

  "showProgress?": "<boolean>"
  // Se true, mostra la barra di progressione con i frammenti
  // scoperti e quelli ancora nascosti.
}
\end{lstlisting}


\section{Contesto di esecuzione (\texttt{CommandContext})}

Per completezza, si riporta la struttura del contesto passato a ogni esecuzione di comando. Questa interfaccia non è configurabile dall'autore, ma è costruita automaticamente dal sistema a runtime.

\begin{lstlisting}[language=json]
{
  "command": "<string>",            // Nome del comando eseguito
  "args": ["<string>"],             // Argomenti passati
  "currentPath": "<string>",        // Directory corrente
  "filesystem": { /* FilesystemStructure */ },
  "environment": { /* Record<string, string> */ },
  "bootStage": "<string>",          // Stadio di boot corrente
  "discoveredFlags": ["<string>"],  // Flag gia' scoperte
  "history": ["<string>"],          // Storico dei comandi eseguiti
  "state": { /* Record<string, any> */ },  // Variabili di stato
  "tab": "<string>"                 // ID del tab corrente
}
\end{lstlisting}


\section{Risultato dell'esecuzione (\texttt{CommandExecutionResult})}

Analogamente, la struttura del risultato restituito dall'esecutore al termine di ogni comando:

\begin{lstlisting}[language=json]
{
  "success": "<boolean>",           // Esecuzione riuscita
  "output": ["<string>"],           // Righe di output
  "error?": "<string>",             // Messaggio di errore (se fallita)
  "sideEffects?": {
    "newFlags?": ["<string>"],      // Flag sbloccate dal comando
    "stateChanges?": {},            // Variabili di stato modificate
    "stageChange?": "<string>",     // Nuovo stadio di boot
    "pathChange?": "<string>"       // Nuova directory corrente
  }
}
\end{lstlisting}


\section{Esempio: configurazione del tab UART Console}

Si riporta un estratto semplificato della configurazione del tab ``UART Console'' dal caso di studio, che illustra l'uso combinato di comandi built-in, comandi custom con vincoli di boot stage, output lookup con sblocco condizionale di flag, e sequenza di boot multi-stadio.

\begin{lstlisting}[language=json]
{
  "id": "uart",
  "name": "UART Console",
  "initialPath": "/",
  "filesystem": {
    "tree": {
      "etc": {
        "passwd": "root:x:0:0:root:/root:/bin/sh",
        "shadow": "root:$1$GTN.gpri$DlSyKvZKMR9A9Uj9e9wR3/:15771:0:99999:7:::"
      },
      "proc": {
        "cpuinfo": "system type  : QCA953x\nprocessor  : 0\ncpu model  : MIPS 24Kc V7.4"
      }
    },
    "directories": {},
    "files": {},
    "fileDefaults": { "permissions": "-rw-r--r--", "owner": "root" }
  },
  "commands": {
    "ls":  { "handler": "builtin", "builtinType": "ls",
             "description": "List directory contents" },
    "cat": { "handler": "builtin", "builtinType": "cat",
             "description": "Display file contents",
             "constraints": { "arguments": { "min": 1 } } },
    "printenv": {
      "description": "Print environment variables",
      "handler": "custom",
      "constraints": { "state": { "bootStage": "uboot_shell" } },
      "output": {
        "type": "static",
        "lines": [
          "bootargs=console=ttyS0,115200 root=31:2 rootfstype=jffs2 init=/sbin/init",
          "bootcmd=bootm 0x9f020000",
          "bootdelay=1",
          "baudrate=115200"
        ]
      },
      "sideEffects": { "unlockFlags": ["boot"] }
    },
    "strings": {
      "description": "Print printable characters in files",
      "handler": "custom",
      "constraints": { "arguments": { "min": 1 } },
      "output": {
        "type": "lookup",
        "argIndex": 0,
        "matchType": "contains",
        "table": {
          "mtdblock3": ["TP-LINK_807C", "WPA2-PSK",
                        "MyW1F1P@ssw0rd!", "192.168.0.1"],
          "httpd":     ["execFormatCmd", "system", "popen"],
          "backdoorTest": ["socket", "connect", "/bin/sh",
                           "192.168.0.100", "4444"]
        }
      },
      "sideEffects": {
        "unlockFlags": [
          { "id": "leak",
            "condition": { "type": "argument", "index": 0,
                           "contains": "mtdblock3" } },
          { "id": "inject",
            "condition": { "type": "argument", "index": 0,
                           "contains": "httpd" } },
          { "id": "shell",
            "condition": { "type": "argument", "index": 0,
                           "contains": "backdoorTest" } }
        ]
      }
    },
    "root": {
      "description": "Login as root",
      "handler": "custom",
      "constraints": { "state": { "bootStage": "login" } },
      "sideEffects": {
        "setState": { "bootStage": "password", "loginUser": "root" }
      }
    },
    "sohoadmin": {
      "description": "Enter password",
      "handler": "custom",
      "constraints": { "state": { "bootStage": "password" } },
      "output": {
        "type": "static",
        "lines": ["", "BusyBox v1.01 Built-in shell (ash)",
                  "Enter 'help' for a list of built-in commands.", ""]
      },
      "sideEffects": {
        "unlockFlags": ["root"],
        "setState": { "bootStage": "shell", "currentPath": "/" }
      }
    }
  },
  "bootSequence": {
    "initialStage": "connecting",
    "stages": [
      { "id": "connecting", "name": "Connecting",
        "lines": ["Connecting to UART interface..."],
        "duration": 1500, "nextStage": "booting" },
      { "id": "booting", "name": "Booting",
        "lines": ["U-Boot 1.1.4 (Jun 16 2015)",
                  "DRAM: 32 MB", "Flash: 4 MB"],
        "duration": 2000, "nextStage": "uboot_wait" },
      { "id": "uboot_wait", "name": "U-Boot Wait",
        "lines": ["Hit any key to stop autoboot: 3"],
        "duration": 1000, "nextStage": "uboot_shell",
        "prompt": "Press Enter to access U-Boot console" },
      { "id": "uboot_shell", "name": "U-Boot Shell",
        "lines": ["U-Boot> "], "prompt": "U-Boot> " },
      { "id": "kernel_boot", "name": "Kernel Boot",
        "lines": ["Booting QCA953x", "Linux version 2.6.31",
                  "VFS: Mounted root (squashfs) readonly"],
        "duration": 2500, "nextStage": "login" },
      { "id": "login", "name": "Login",
        "lines": ["TP-LINK login: "], "prompt": "Login: " },
      { "id": "password", "name": "Password",
        "lines": [], "prompt": "Password: " },
      { "id": "shell", "name": "Shell",
        "lines": ["BusyBox v1.01 Built-in shell (ash)"],
        "prompt": "# " }
    ]
  },
  "environment": {
    "USER": "root", "HOME": "/root",
    "PATH": "/bin:/sbin:/usr/bin:/usr/sbin", "SHELL": "/bin/sh"
  },
  "defaultConstraints": { "state": { "bootStage": "shell" } }
}
\end{lstlisting}

Questo esempio illustra i principali pattern di configurazione:

\begin{itemize}
    \item I comandi \texttt{ls} e \texttt{cat} utilizzano handler built-in con vincoli sugli argomenti;
    \item Il comando \texttt{printenv} è limitato allo stadio \texttt{uboot\_shell} e sblocca incondizionatamente la flag \texttt{boot};
    \item Il comando \texttt{strings} usa un output di tipo lookup con sblocco condizionale: la flag sbloccata dipende dall'argomento passato dallo studente;
    \item I comandi \texttt{root} e \texttt{sohoadmin} gestiscono il flusso di login tramite transizioni di stato (\texttt{bootStage});
    \item I \texttt{defaultConstraints} del tab limitano tutti i comandi (tranne quelli con vincoli propri) allo stadio \texttt{shell}, impedendo l'uso di comandi filesystem durante il prompt U-Boot.
\end{itemize}


\end{document}
